\documentclass[UTF8]{ctexart}
\usepackage{xeCJK}
\usepackage{geometry}
\geometry{left=2cm,right=2cm,top=2.5cm,bottom=2.5cm}
\setlength{\headheight}{12.7pt}

\usepackage{titlesec}
\titleformat{\section}
  {\normalfont\large\bfseries}
  {\thesection}
  {1em}
  {}
\titlespacing*{\section}
  {0pt}
  {3.5ex plus 1ex minus .2ex}
  {2.3ex plus .2ex}

\usepackage{amsmath}
\usepackage{amssymb}
\usepackage{amsthm}
\usepackage{enumitem}
\setlist[enumerate]{itemsep=0pt,topsep=0pt,parsep=0pt,leftmargin=0pt,itemindent=2.5em}
\setlist[itemize]{itemsep=0pt,topsep=0pt,parsep=0pt,leftmargin=0pt,itemindent=2.5em}
\usepackage{fancyhdr}
\pagestyle{fancy}
\fancyhf{}
\usepackage{graphicx}
\usepackage{hyperref}
\usepackage{indentfirst}
\usepackage{lmodern}


\begin{document}
\setlength{\abovedisplayskip}{1pt}
\setlength{\belowdisplayskip}{1pt}

\section{部分正确性公式}

\hypertarget{sec:定义1.1}{}
\noindent\textbf{定义1.1 (部分正确性公式)}

对任意的断言(即公式)$P,Q$和while程序$S$, 称
\[
    \{P\}\,S\,\{Q\}
\]
为一个\textbf{部分正确性公式}.

\vspace{10pt}

\hypertarget{sec:定义1.2}{}
\noindent\textbf{定义1.2 (部分正确性公式的语义)}

给定断言$P,Q$和程序$S$, 如果$\mathcal{M}[\![S]\!](P)\subset Q$, 即
\[
    \forall\sigma,\tau\in\Sigma:\quad
    \,\sigma\models P\land [\![S]\!](\sigma)=\tau\quad\to\quad\tau\models Q\,,
\]
就记$\models\{P\}\,S\,\{Q\}$.

\vspace{10pt}

\hypertarget{sec:定义1.3}{}
\noindent\textbf{定义1.3 (PW证明系统)}

while程序部分正确性公式的PW(partial, while)证明系统:
\begin{enumerate}
    \item \textbf{skip公理}:
    \[
        \{P\}\,\mathrm{skip}\,\{P\},
    \]
    \item \textbf{assignment公理}(后向形式): 当$P[u:=t]$是自由替换时,
    \[
        \{\,P[u:=t]\,\}\,u:=t\,\{P\},
    \]
    \item \textbf{composition规则}:
    \[
        \frac{\{P\}\,S_1\,\{R\},\,\{R\}\,S_2\,\{Q\}}
        {\{P\}\,S_1;S_2\,\{Q\}},
    \]
    \item \textbf{conditional规则}:
    \[
        \frac{\{P\land B\}\,S_1\,\{Q\},\,\{P\land\neg\,B\}\,S_2\,\{Q\}}
        {\{P\}\,\mathbf{if}\,\,B\,\,\mathbf{then}\,\,S_1\,\,\mathbf{else}
        \,\,S_2\,\,\mathbf{fi}\{Q\}},
    \]
    \item \textbf{loop规则}:
    \[
        \frac{\{P\land B\}\,S\,\{P\}}
        {\{P\}\,\mathbf{while}\,\,B\,\,\mathbf{do}\,\,S\,\,\mathbf{od}\,
        \{P\land\neg\,B\}},
    \]
    \item \textbf{consequence规则}:
    \[
        \frac{P\to P_1,\,\{P_1\}\,S\,\{Q_1\},\,Q_1\to Q}
        {\{P\}\,S\,\{Q\}}.
    \]
\end{enumerate}

\vspace{15pt}

assignment公理有前向形式: 当与$u$同类型的变元$v$不出现在$u,t,P$中时,
\[
    \vdash_{\mathrm{PW}}
    \{P\}\,u:=t\,\{\,\exists v\,(\,u\equiv t[u:=v]\land P[u:=v]\,)\,\}.
\]
它可由后向形式推出: 记$Q=\exists v\,(\,u\equiv t[u:=v]\land P[u:=v]\,)$, 则易证
\begin{enumerate}
    \item $P\to t\equiv t[u:=u]\land P[u:=u]\to
    \exists v\,(\,t\equiv t[u:=v]\land P[u:=v]\,)$,
    \item $Q[u:=t]=\exists v\,(\,t\equiv t[u:=v]\land P[u:=v]\,)$,
    并且是自由替换,
\end{enumerate}
所以$P\to Q[u:=t]$且$\vdash_{\mathrm{PW}}\{Q[u:=t]\}\,u:=t\,\{Q\}$,
依consequence规则得$\vdash_{\mathrm{PW}}\{P\}\,u:=t\,\{Q\}$.





\newpage

\section{PW的可靠性}

\hypertarget{sec:定理2.1}{}
\noindent\textbf{定理2.1 (PW证明系统的可靠性)}

对任意的断言$P,Q$和程序$S$, 如果$\vdash_{\mathrm{PW}}\{P\}\,S\,\{Q\}$,
那么$\models\{P\}\,S\,\{Q\}$.

\noindent{\kaishu 证明.}

对PW规则归纳证明.

\begin{enumerate}
    \item $\mathcal{M}[\![\mathrm{skip}]\!](P)=P$, 所以$\models
    \{P\}\,\mathrm{skip}\,\{P\}$,

    \item 当$P[u:=t]$是自由替换时, 如果$\sigma\models P[u:=t]$
    且$[\![u:=t]\!](\sigma)=\tau$, 那么根据替换定理有
    \[
        \tau=\sigma[u:=\sigma(t)]\models P,
    \]
    所以$\models\{\,P[u:=t]\,\}\,u:=t\,\{P\}$.

    \item 假设$\models\{P\}\,S_1\,\{R\}$, $\models\{R\}\,S_2\,\{Q\}$, 则
    \[
        \mathcal{M}[\![S_1;S_2]\!](P)=\mathcal{M}[\![S_2]\!]
        \mathcal{M}[\![S_1]\!](P)\subset\mathcal{M}[\![S_2]\!](R)\subset Q,
    \]
    所以$\models\{P\}\,S_1;S_2\,\{Q\}$.

    \item 假设$\models\{P\land B\}\,S_1\,\{Q\}$,
    $\models\{P\land\neg\,B\}\,S_2\,\{Q\}$, 则
    \[
        \mathcal{M}[\![\,\mathbf{if}\,\,B\,\,\mathbf{then}\,\,S_1\,\,
        \mathbf{else}\,\,S_2\,\,\mathbf{fi}\,]\!](P)=
        \mathcal{M}[\![S_1]\!](P\land B)\cup
        \mathcal{M}[\![S_2]\!](P\land\neg B)\subset Q,
    \]
    所以$\models\{P\}\,\mathbf{if}\,\,B\,\,\mathbf{then}\,\,S_1\,\,
    \mathbf{else}\,\,S_2\,\,\mathbf{fi}\,\{Q\}$.

    \item 假设$\models\{P\land B\}\,S\,\{P\}$,
    对$k$归纳证明$\mathcal{M}[\![A^k]\!](P)\subset P\land\neg\,B$.

    \hspace{2.17em}
    首先$\mathcal{M}[\![A^0]\!](P)=\varnothing$; 假设命题对$k$成立, 那么
    \vspace{3pt}
    \[\begin{aligned}
        \mathcal{M}[\![A^{k+1}]\!](P)=\,&
        \mathcal{M}[\![\,\mathbf{if}\,\,B\,\,\mathbf{then}\,\,S;A^k\,\,
        \mathbf{else}\,\,\mathrm{skip}\,\,\mathbf{fi}\,]\!](P)\\
        =\,&\mathcal{M}[\![S;A^k]\!](P\land B)\cup
        \mathcal{M}[\![\mathrm{skip}]\!](P\land\neg\,B)\\
        =\,&\mathcal{M}[\![A^k]\!]\mathcal{M}[\![S]\!](P\land B)
        \cup(P\land\neg\,B)\\
        \subset\,&\mathcal{M}[\![A^k]\!](P)\cup(P\land\neg\,B)
        =P\land\neg\,B,\\[3pt]
    \end{aligned}\]
    归纳成立, 所以
    \[
        \mathcal{M}[\![\,\mathbf{while}\,\,B\,\,\mathbf{do}\,\,
        S\,\,\mathbf{od}\,]\!](P)=\bigcup_{k\in\mathbb{N}}\mathcal{M}
        [\![A^k]\!](P)\subset P\land\neg\,B,\\[-5pt]
    \]
    即$\models\{P\}\,\mathbf{while}\,\,B\,\,\mathbf{do}\,\,S
    \,\,\mathbf{od}\,\{P\land\neg\,B\}$.
    
    \item 假设$P\to P_1$, $\models\{P_1\}\,S\,\{Q_1\}$且$Q\to Q_1$, 则
    \[
        \mathcal{M}[\![S]\!](P)\subset\mathcal{M}[\![S]\!](P_1)
        \subset Q_1\subset Q,
    \]
    即$\models\{P\}\,S\,\{Q\}$. \hfill $\qedsymbol$
\end{enumerate}

\end{document}
